%% Supplementary Material for the SC-STEM paper (Metron special issue).
%% Structure follows the first (arXiv) version of the SC-FH paper: the
%% derivations and the extended per-design results live here, the manuscript
%% refers to them by section number. Cross-document \ref is not available, so
%% references to the manuscript are textual.

\documentclass[11pt]{article}

\usepackage[margin=1in]{geometry}
\usepackage{setspace}
\usepackage{amsmath, amssymb, mathtools}
\usepackage{graphicx}
\usepackage{booktabs}
\usepackage{adjustbox}
\usepackage{microtype}
\usepackage{xcolor}
\usepackage{enumitem}
\usepackage{algorithm}
\usepackage{algpseudocode}
\usepackage{hyperref}
\usepackage[round]{natbib}

\hypersetup{colorlinks=true, linkcolor=blue, citecolor=blue, urlcolor=blue}
\graphicspath{{Figures/}}

\newcommand{\TODO}[1]{\textcolor{red}{\textbf{[TO RUN:} #1\textbf{]}}}
\newcommand{\Stem}{\texttt{Stem}}
\newcommand{\Rlang}{\textsf{R}}
\DeclareMathOperator*{\argmax}{arg\,max}

%% Supplementary numbering: sections S1, S2, ...; figures S1, ...; tables S1, ...
\renewcommand{\thesection}{S\arabic{section}}
\renewcommand{\thefigure}{S\arabic{figure}}
\renewcommand{\thetable}{S\arabic{table}}
\renewcommand{\theequation}{S\arabic{equation}}

\title{Supplementary Material for\\
       ``Spatio-Temporal Models under Spatial Regimes: the Spatially-Clustered
       STEM Family and its Implementation in the \Stem{} \Rlang{} Package''}
\author{Paolo Maranzano}
\date{\today}

\begin{document}
\maketitle

\noindent This document collects the derivations and the extended results accompanying the main paper. Section~\ref{sm:missing} derives the completed sufficient statistics that the EM algorithm requires when the response is incompletely observed, which the manuscript states without proof. Section~\ref{sm:ic} derives the behaviour of the information criteria as the dimensions grow. Sections~\ref{sm:lambda} and~\ref{sm:order} document the ridge as the package implements it: the exact scale of $\lambda$, the switches that reduce the model to penalized linear regression, the relation between one $\lambda$ and the regimes, and the order in which the hyperparameters have to be chosen. Section~\ref{sm:designs} details the simulation designs, and Sections~\ref{sm:scenarios}--\ref{sm:balance} report the full per-design results. Section~\ref{sm:boot} validates the parametric bootstrap and Section~\ref{sm:app} completes the empirical application. Section~\ref{sm:code} documents the reproduction code. All figures and tables are generated by the analysis scripts.

%%%%%%%%%%%%%%%%%%%%%%%%%%%%%%%%%%%%%%%%%%%%%%%%%%%%%%%%%%%%%%%%%%%%%%%%%%%%%%%
\section{Missing responses: the completed sufficient statistics}\label{sm:missing}
%%%%%%%%%%%%%%%%%%%%%%%%%%%%%%%%%%%%%%%%%%%%%%%%%%%%%%%%%%%%%%%%%%%%%%%%%%%%%%%

Section~4.5 of the manuscript states the form of the expected complete-data sufficient statistic of the measurement equation when part of the response is unobserved. This section derives it. Throughout, the regime index $k$ is fixed and suppressed, $Z$ denotes the $d \times p$ loading matrix, $\Sigma_e$ the measurement covariance, and $\mathcal{Z}$ the whole of the observed data.

\subsection{Setting}

The measurement equation is $z_t = X_t \beta + Z y_t + e_t$ with $e_t \sim N(0, \Sigma_e)$, independent across $t$ and of the latent process. Partition the locations at time $t$ into the observed set $O$ and the missing set $M$, and write the corresponding blocks of any vector or matrix with those subscripts. The EM algorithm treats both the latent states \emph{and} the unobserved responses as missing data, so the E-step must return
\begin{equation}\label{eq:sm:target}
  \mathrm{E}\!\left[e_t e_t^\top \,\middle|\, \mathcal{Z}\right],
  \qquad\text{and}\qquad
  \mathrm{E}\!\left[z_t \,\middle|\, \mathcal{Z}\right],
\end{equation}
the first entering the updates of the variance components and of the spatial covariance parameters, the second the update of $\beta$.

The Kalman smoother, run on the selected measurement equation of \citet[Sect.~4.10]{DurbinKoopman2012}, supplies
\begin{equation}\label{eq:sm:smooth}
  \hat{y}_t = \mathrm{E}\!\left[y_t \mid \mathcal{Z}\right],
  \qquad
  C^s_t = \mathrm{Var}\!\left[y_t \mid \mathcal{Z}\right].
\end{equation}

\subsection{The conditional law of the measurement error}

Conditionally on the latent state $y_t$, the measurement error is $e_t = z_t - X_t\beta - Z y_t$, so its observed block $e_{O,t}$ is a known linear function of $y_t$ once $z_{O,t}$ is observed, whereas its missing block is not observed at all. Because $e_t$ is Gaussian with covariance $\Sigma_e$, the standard conditioning formula gives
\begin{equation}\label{eq:sm:cond}
  e_{M,t} \mid e_{O,t} \;\sim\; N\!\left(P\, e_{O,t},\; \Omega\right),
  \qquad
  P = \left[\Sigma_e\right]_{MO}\left(\left[\Sigma_e\right]_{OO}\right)^{-1},
  \qquad
  \Omega = \left[\Sigma_e\right]_{MM} - P \left[\Sigma_e\right]_{OM},
\end{equation}
and this conditional law does not depend on $y_t$: the state enters $e_{O,t}$ but not the regression of $e_{M,t}$ on it. This is the point at which the spatio-temporal case departs from the textbook treatment. \citet{DurbinKoopman2012} note that a missing element of the observation vector can be estimated by the corresponding element of $Z_t \hat{y}_t$; that is exact when $\Sigma_e$ is diagonal, in which case $P = 0$, and it is not exact here, where the exponential spatial covariance couples the locations and $P \neq 0$.

Introduce the $d \times |O|$ matrix $S$ that lifts a vector defined on the observed rows to the full vector,
\begin{equation}\label{eq:sm:S}
  S_{O\cdot} = I_{|O|}, \qquad S_{M\cdot} = P,
\end{equation}
and let $\tilde{\Omega}$ be the $d \times d$ matrix with $\Omega$ in the $(M,M)$ block and zeros elsewhere. Then \eqref{eq:sm:cond} can be written compactly as
\begin{equation}\label{eq:sm:compact}
  e_t \mid e_{O,t} \;\sim\; N\!\left(S\, e_{O,t},\; \tilde{\Omega}\right).
\end{equation}

\subsection{The completed second moment}

Write $r_{O,t} = z_{O,t} - X_{O,t}\beta - Z_{O} \hat{y}_t$ for the residual on the observed rows, evaluated at the smoothed state. Applying the tower property with respect to $y_t \mid \mathcal{Z}$,
\begin{align}
  \mathrm{E}\!\left[e_t \mid \mathcal{Z}\right]
    &= \mathrm{E}\!\left[\, \mathrm{E}\!\left[e_t \mid y_t, \mathcal{Z}\right] \,\middle|\, \mathcal{Z}\right]
     = S\, r_{O,t}, \label{eq:sm:mean}\\[2pt]
  \mathrm{Var}\!\left[e_t \mid \mathcal{Z}\right]
    &= \underbrace{\mathrm{E}\!\left[\mathrm{Var}\!\left(e_t \mid y_t, \mathcal{Z}\right) \middle| \mathcal{Z}\right]}_{= \;\tilde{\Omega}}
     + \underbrace{\mathrm{Var}\!\left[\mathrm{E}\!\left(e_t \mid y_t, \mathcal{Z}\right) \middle| \mathcal{Z}\right]}_{= \; S Z_O C^s_t Z_O^\top S^\top}, \label{eq:sm:var}
\end{align}
where the second term follows because, given $y_t$, the conditional mean of $e_t$ is $S\left(z_{O,t} - X_{O,t}\beta - Z_O y_t\right)$, an affine function of $y_t$ with matrix $-S Z_O$, whose variance under the smoothing distribution is $S Z_O C^s_t Z_O^\top S^\top$. Combining \eqref{eq:sm:mean} and \eqref{eq:sm:var},
\begin{equation}\label{eq:sm:BB}
  \boxed{\;
  \mathrm{E}\!\left[e_t e_t^\top \,\middle|\, \mathcal{Z}\right]
  \;=\; S\, Z_O\, C^s_t\, Z_O^\top S^\top \;+\; \tilde{\Omega} \;+\; \left(S r_{O,t}\right)\left(S r_{O,t}\right)^{\top}. \;}
\end{equation}

Two degenerate cases confirm the expression. If nothing is missing, then $M = \emptyset$, $S = I_d$ and $\tilde{\Omega} = 0$, and \eqref{eq:sm:BB} reduces to $Z C^s_t Z^\top + r_t r_t^\top$, the complete-data statistic. If nothing is observed, then $O = \emptyset$, $S$ has no columns and $\tilde{\Omega} = \Sigma_e$, so \eqref{eq:sm:BB} returns $\Sigma_e$: with no information at that time point the measurement error keeps its unconditional variance, and contributes nothing else.

\subsection{The completed observation and the update of \texorpdfstring{$\beta$}{beta}}

From \eqref{eq:sm:mean}, the conditional expectation of the observation vector is
\begin{equation}\label{eq:sm:zhat}
  \hat{z}_t = \mathrm{E}\!\left[z_t \mid \mathcal{Z}\right]
            = X_t\beta + Z\hat{y}_t + S\, r_{O,t},
\end{equation}
which equals $z_t$ on the observed rows, as it must, and equals $X_{M,t}\beta + Z_M \hat{y}_t + P r_{O,t}$ on the missing ones. Differentiating the expected complete-data log-likelihood with respect to $\beta$ and setting the derivative to zero gives
\begin{equation}\label{eq:sm:beta}
  \hat{\beta} = \left(\sum_{t=1}^{T} X_t^\top \Sigma_e^{-1} X_t\right)^{-1}
                \sum_{t=1}^{T} X_t^\top \Sigma_e^{-1}\left(\hat{z}_t - Z \hat{y}_t\right),
\end{equation}
so the only change with respect to the complete-data update is the replacement of $z_t$ by $\hat{z}_t$: the design matrix and the covariance keep their full dimension. For the same reason the divisor of the variance update remains the complete-data count $dT$ and not the number of observed values, since \eqref{eq:sm:BB} already carries the conditional variance of whatever was not observed.

\subsection{Cost}

The matrices $P$ and $\Omega$ depend on the time point only through its missingness pattern, so they are computed once per distinct pattern and cached. In a monitoring network a station is typically out of service for a stretch of consecutive days, so the number of distinct patterns is far smaller than $T$, and the additional cost of the treatment is negligible relative to the filtering pass.

%%%%%%%%%%%%%%%%%%%%%%%%%%%%%%%%%%%%%%%%%%%%%%%%%%%%%%%%%%%%%%%%%%%%%%%%%%%%%%%
\section{Information criteria as \texorpdfstring{$T$}{T} and \texorpdfstring{$d$}{d} grow}\label{sm:ic}
%%%%%%%%%%%%%%%%%%%%%%%%%%%%%%%%%%%%%%%%%%%%%%%%%%%%%%%%%%%%%%%%%%%%%%%%%%%%%%%

This section derives what the manuscript states in words: how the likelihood-based criteria behave when the series lengthen and when the network widens, and why the two directions are not symmetric.

\subsection{Notation and the two quantities to compare}

Write $q = r + 3 + 3p$ for the number of parameters of one regime, so that a $k$-regime configuration has $\mathcal{K}(k) = kq$ continuous parameters, and $n = dT$ for the number of observations entering the likelihood. Let $\hat\ell(k)$ be the maximized log-likelihood at $k$ regimes, the maximum being taken over the parameters \emph{and} over the partitions. A criterion moves from $k$ to $k+1$ when
\begin{equation}\label{eq:sm:rule}
  2\,\Delta\hat\ell(k) \;>\; \Delta\mathrm{pen}(k),
  \qquad
  \Delta\hat\ell(k) = \hat\ell(k+1) - \hat\ell(k).
\end{equation}
Everything follows from the orders of the two sides.

\subsection{The likelihood gain}

\paragraph{Under the alternative.} If the truth has $k_0 > k$ regimes, the best $k$-regime approximation sits at a strictly positive Kullback--Leibler distance $\delta_k$ per observation from the truth, and
\begin{equation}\label{eq:sm:gainalt}
  2\,\Delta\hat\ell(k) \;\asymp\; 2\, n\, \delta_k \;=\; O_p(dT).
\end{equation}
The gain is proportional to the number of observations the extra regime describes: it grows in \emph{both} $d$ and $T$.

\paragraph{Under the null.} If the truth has $k$ regimes and we fit $k+1$, the extra regime is superfluous. For a \emph{fixed} refinement of the true partition the likelihood ratio is asymptotically $\chi^2_q$, hence $O_p(1)$ in $T$: lengthening the series does not help a spurious regime. What does help it is the search. The maximum is taken over the partitions of $d$ locations into $k+1$ groups, a set of cardinality $S(d, k+1)$, the Stirling number of the second kind, and $\log S(d,k) = d\log k - \log k! + o(d)$. The standard extreme-value heuristic for the maximum of many $\chi^2_q$ variables --- the same heuristic that underlies the extended BIC --- gives
\begin{equation}\label{eq:sm:gainnull}
  2\,\Delta\hat\ell(k) \;\asymp\; 2\left[\log S(d,k+1) - \log S(d,k)\right]
  \;\approx\; 2\,d \log\!\left(\tfrac{k+1}{k}\right) \;=\; O_p(d).
\end{equation}
This is the first asymmetry, and it is the whole story in miniature: \textbf{under the null the spurious gain grows linearly in $d$ and not at all in $T$.} The variables being maximized over are strongly dependent, since neighbouring partitions share most of their labels, so \eqref{eq:sm:gainnull} is an upper bound on the order rather than an equality; it is the conservative and the relevant one.

\subsection{The penalty increments}

Writing $\mathcal{K} = kq$ and taking the increment from $k$ to $k+1$,
\begin{align}
  \Delta\mathrm{AIC}  &= 2q, \label{eq:sm:aic}\\
  \Delta\mathrm{AICc} &= 2q + \frac{2q^2(2k+1)}{n} + O(n^{-2}), \label{eq:sm:aicc}\\
  \Delta\mathrm{BIC}  &= q \log(dT) = q(\log d + \log T), \label{eq:sm:bic}\\
  \Delta\mathrm{EBIC}_\gamma &= q \log(dT) + 2\gamma\left[\log S(d,k+1) - \log S(d,k)\right]
     \;\approx\; q\log(dT) + 2\gamma\, d \log\!\left(\tfrac{k+1}{k}\right). \label{eq:sm:ebic}
\end{align}

\subsection{Consistency, direction by direction}

Combining \eqref{eq:sm:gainnull} with \eqref{eq:sm:aic}--\eqref{eq:sm:ebic}, a criterion avoids a spurious regime when its increment exceeds $c\,d$ with $c = 2\log\{(k+1)/k\}$.

\paragraph{AIC and AICc.} The increment is $2q$, a constant. The condition $2q > c\,d$ holds only for $d < 2q/c$, so \textbf{both are inconsistent, and the failure worsens linearly in $d$}. Neither depends on $T$ at all. The corrected version adds nothing: by \eqref{eq:sm:aicc} the correction is $O(\mathcal{K}^2/n) = O(k^2q^2/(dT))$ and vanishes in \emph{both} directions. The AICc is built for the regime in which $\mathcal{K}/n$ is not negligible; in a spatio-temporal panel $\mathcal{K}/n = kq/(dT) \to 0$ however the sample grows, because $\mathcal{K}$ is fixed by the specification and $n$ multiplies $d$ by $T$.

\paragraph{BIC.} The condition is $q(\log d + \log T) > c\,d$. Hence
\begin{itemize}[leftmargin=1.4em,itemsep=2pt,topsep=3pt]
  \item \emph{$T \to \infty$ with $d$ fixed}: the left side diverges, the right is fixed. \textbf{The BIC is consistent in $T$.}
  \item \emph{$d \to \infty$ with $T$ fixed}: the left side grows like $q\log d$, the right like $c\,d$. \textbf{The BIC over-selects, and does so for every large enough network.}
  \item \emph{both growing}: consistency requires $d / \log(dT) \to 0$, that is $d = o(\log T)$ --- the network must grow only logarithmically in the length of the series.
\end{itemize}
Lengthening the record therefore helps the BIC and widening the network hurts it, which is the opposite of the intuition that more data is more information.

\paragraph{Extended BIC.} Here the penalty carries a term of the same order as the null gain. The condition becomes
\[
  q\log(dT) + 2\gamma\, d \log\!\left(\tfrac{k+1}{k}\right) \;>\; 2\,d\log\!\left(\tfrac{k+1}{k}\right),
\]
which holds for all large $d$ as soon as $\gamma \ge 1$. \textbf{The extended BIC is the only member of the family whose penalty has the right order in $d$, and $\gamma = 1$ is exactly the calibration point}, since it charges precisely the multiplicity that \eqref{eq:sm:gainnull} extracts. This is the theoretical counterpart of the numerical finding reported in the manuscript: on the real data the extended BIC does not change the answer, not because it is miscalibrated, but because the observed gain is far above anything the null can produce.

\subsection{Criteria built on the penalized objective}\label{sm:ic:pen}

Everything so far applies to criteria built on the log-likelihood. It is natural to ask what happens if they are built instead on the objective the algorithm actually maximizes,
\begin{equation}\label{eq:sm:Q}
  Q(k,\phi) \;=\; \sum_{i=1}^{d} \ell_i(k_i) \;+\; \beta\, P(\mathcal{P}),
  \qquad \beta = \phi\, c,
\end{equation}
with $P(\mathcal{P})$ the number of neighbouring pairs that share a label. The answer depends on two scalings, and it is worth separating them.

\paragraph{How the Potts term scales.} On an $m$-nearest-neighbour graph over $d$ locations the edge set has $|E| \asymp md/2 = O(d)$ elements, so $P \le |E| = O(d)$. When the regimes are spatially coherent, the pairs that are \emph{not} concordant are those straddling a boundary. For $k$ compact regions in a two-dimensional domain the total interface length is $O(L\sqrt{k})$ with $L$ the linear size of the domain, and the number of nodes within one spacing of it is that length times the linear density $\sqrt{d}/L$. Hence
\begin{equation}\label{eq:sm:cut}
  C(\mathcal{P}) \;\equiv\; |E| - P(\mathcal{P}) \;=\; O\!\left(\sqrt{dk}\right),
  \qquad
  \Delta_k P \;=\; -\,O\!\left(\sqrt{dk}\right).
\end{equation}
The concordant count is therefore of order $d$ but its \emph{increment} in $k$ is only of order $\sqrt{dk}$: a boundary is a lower-dimensional object. At $\phi = 0$ the partition need not be coherent and $C$ can be $O(d)$, which is the one regime where \eqref{eq:sm:cut} does not apply.

\paragraph{How the scale factor scales.} The three options of the implementation differ exactly here. With \texttt{raw} scaling $c = 1$. With \texttt{per-observation} scaling $c = T$. With the default \texttt{auto} scaling $c$ is the median across locations of the spread $\max_k \ell_i(k) - \min_k \ell_i(k)$, divided by the mean degree; each $\ell_i(k)$ is a sum of $T$ terms and the per-period discrepancy between the best and the worst regime is $O(1)$, so
\begin{equation}\label{eq:sm:cscale}
  c \;=\; O(1) \ \text{(raw)}, \qquad c \;=\; O(T) \ \text{(per-observation and auto)} .
\end{equation}

\paragraph{The naive version is not a criterion.} Take $\mathrm{IC}_{\mathrm{pen}} = -2Q + \mathrm{pen}(\mathcal{K}, n)$. For fixed $k$ it is affine and decreasing in $\phi$, with slope $-2cP$, so it is minimized at the largest penalty on the grid whatever the data say: it does not select $\phi$, it merely records that $\phi$ is large. Across $k$ the sign depends on the scaling. With $c$ constant the increment is $2\beta\,\Delta_k C = O(\beta\sqrt{dk})$, a genuine penalty on additional regimes; with the \texttt{auto} rule $c$ is itself computed from the spread of the scores across the $k$ candidate regimes and therefore \emph{grows} with $k$, so $\beta P$ increases with $k$ and the term rewards additional regimes. Under the default scaling the comparison across $k$ is between quantities measured on different scales, which by itself invalidates it.

\paragraph{The missing constant, and what the corrected version is.} The deeper obstruction is that $\beta P$ is the logarithm of a Potts prior on the partition only up to its normalizing constant
\[
  Z(\beta,k) \;=\; \sum_{\text{all } k^d \text{ labelings}} \exp\{\beta P\},
\]
which depends on both quantities being selected. Restoring it changes the term to $\beta P - \log Z$. As $\beta$ grows the partition function is dominated by the $k$ constant labelings, $\log Z = \beta |E| + \log k + o(1)$, so
\begin{equation}\label{eq:sm:normalized}
  \beta P - \log Z \;\longrightarrow\; -\,\beta\, C(\mathcal{P}) \;-\; \log k ,
\end{equation}
and the properly normalized criterion is
\begin{equation}\label{eq:sm:ICnorm}
  \mathrm{IC}_{\mathrm{pen}}^{\,\mathrm{norm}}
  \;=\; \underbrace{-2\hat\ell + \mathrm{pen}(\mathcal{K}, n)}_{\text{the ordinary criterion}}
  \;+\; 2\beta\, C(\mathcal{P}) \;+\; 2\log k .
\end{equation}
The Potts term, once normalized, is neither a reward nor a nuisance: it is a penalty proportional to the number of edges the partition cuts, that is to the length of the boundary it draws. The sign is the opposite of the naive version's, which is why the two select differently.

\paragraph{Orders, and what this buys.} Combining \eqref{eq:sm:cut} and \eqref{eq:sm:cscale}, the increment of the extra term in \eqref{eq:sm:ICnorm} is
\begin{equation}\label{eq:sm:penorder}
  2\beta\,\Delta_k C \;=\;
  \begin{cases}
    O\!\left(\sqrt{d}\right), & c = O(1) \quad \text{(raw)},\\[2pt]
    O\!\left(T\sqrt{d}\right), & c = O(T) \quad \text{(per-observation, auto)} .
  \end{cases}
\end{equation}
Set against the null gain $O(d)$ of \eqref{eq:sm:gainnull} and the alternative gain $O(dT)$ of \eqref{eq:sm:gainalt}:
\begin{itemize}[leftmargin=1.4em,itemsep=2pt,topsep=3pt]
  \item with $c = O(1)$ the extra penalty is $O(\sqrt d)$, below the null gain $O(d)$, so it does not repair the inconsistency in $d$ and is negligible against the alternative;
  \item with $c = O(T)$ it is $O(T\sqrt d)$, which dominates the null gain as soon as $T \gg \sqrt d$ and is smaller than the alternative gain by a factor $\sqrt d$. Such a criterion is therefore consistent in both directions, at the price of a detection handicap of order $\sqrt{d}$ against genuine regimes.
\end{itemize}
So a criterion on the penalized objective can be made to work, but only after restoring $Z$, and only in the second scaling. Two caveats remain, and they are decisive in practice. Evaluating $Z(\beta,k)$ requires path sampling at every configuration of the grid, which costs more than the cross-validation it would replace. And \eqref{eq:sm:ICnorm} is monotone in $\phi$ in the direction of coherence: raising $\phi$ raises the charge on every cut edge and ultimately drives the selection to $k = 1$, which has none. The criterion can rank $k$ \emph{given} $\phi$; it cannot choose $\phi$, which stays what it was, the strength of a prior fixed from outside the data.

\subsection{Two effective sample sizes, and where the AICc diverges}

The tension is that the model carries parameters of two kinds. The regime parameters $\Psi_k$ are continuous and their information accrues at rate $d_k T$. The partition is discrete, has one component per location, and each label is identified by the $T$ observations of that location: its dimension grows with $d$ while its information per component grows with $T$.

If one insists on reading the partition as the object being selected, the relevant sample size is $d$ and not $dT$, and the AICc correction becomes $2\mathcal{K}(\mathcal{K}+1)/(d - \mathcal{K} - 1)$. This diverges, and then changes sign, when
\begin{equation}\label{eq:sm:bound}
  \mathcal{K} = kq \;\ge\; d - 1
  \qquad\Longleftrightarrow\qquad
  k \;\ge\; \frac{d-1}{q},
\end{equation}
which for the application of the manuscript, $d = 36$ and $q = 9$, gives $k \ge 3.9$. The criterion breaks down exactly at $k = 4$. This is not a numerical accident: it is the same bound, reached from the other side, that the minimum-size constraint of the estimation algorithm imposes, since a regime needs at least $r+2$ locations and hence $k \le d/(r+2)$. A network of $36$ stations cannot identify $36$ regime-specific parameters, whichever way one counts.

\subsection{Dependence, and the ambiguity of \texorpdfstring{$n$}{n}}

A last remark sharpens the point. All of the above uses $n = dT$, the number of \emph{observations}. Under dependence that overstates the information. For a latent process with persistence $G$, the variance of a temporal average over $T$ points is inflated by $(1+G)/(1-G)$, so an effective count is
\begin{equation}\label{eq:sm:neff}
  T_{\mathrm{eff}} \;\approx\; T\,\frac{1-G}{1+G},
\end{equation}
which at the value $\hat{G} = 0.97$ estimated on the real data equals $T \times 0.015$: the $1{,}826$ days carry the information of roughly $28$ independent ones. Substituting $n_{\mathrm{eff}}$ for $n$ in \eqref{eq:sm:bic} lowers the BIC increment further, from $q\log(dT)$ to $q\log(d\,T_{\mathrm{eff}})$, and makes the criterion weaker still.

The conclusion is not that one of these choices of $n$ is correct and the others are wrong. It is that the answer moves by a large factor depending on a quantity the criteria do not determine, while the out-of-sample error of Section~4.6 of the manuscript is defined without reference to it. That, and not a numerical accident, is why the selection is settled by cross-validation.

\subsection{Why no in-sample criterion can settle the question}

There is a final, structural reason, independent of the rates above. Consistency arguments require an asymptotic frame, and the two available frames answer different questions.

With $T \to \infty$ and $d$ fixed, every label is consistently estimated, since each is identified by an unboundedly informative series, and the BIC is consistent by the argument above. But $d$ is fixed: the model is fitted to a closed set of locations and there is no sense in which its partition extends beyond them.

With $d \to \infty$, which is the frame in which spatial prediction is meaningful, the number of labels grows with the sample. The partition is then an \emph{incidental parameter} in the sense of \citet{NeymanScott1948}: its dimension is proportional to the sample size in the spatial direction, and it is not consistently estimated at fixed $T$. No penalty that is a function of $\mathcal{K}$ and $n$ alone controls the error made in assigning a location to a regime, because that error does not vanish in this frame.

This is why a criterion that measures out-of-sample error, and in particular one whose folds withhold whole locations, answers a question the in-sample criteria cannot even pose.

%%%%%%%%%%%%%%%%%%%%%%%%%%%%%%%%%%%%%%%%%%%%%%%%%%%%%%%%%%%%%%%%%%%%%%%%%%%%%%%
\section{The ridge in the implementation: three choices, stated exactly}\label{sm:lambda}
%%%%%%%%%%%%%%%%%%%%%%%%%%%%%%%%%%%%%%%%%%%%%%%%%%%%%%%%%%%%%%%%%%%%%%%%%%%%%%%

Section~2.7 of the manuscript writes the ridge on the original scale of the coefficients. The implementation departs from that notation in three places, each of which changes what a given value of $\lambda$ means. This section states them exactly, because a reader who reproduces the estimator from \eqref{eq:sm:ridge-orig} below alone will not obtain the numbers the package returns.

\subsection{The scale of \texorpdfstring{$\lambda$}{lambda}, and why the ridge needs no further one}

\paragraph{The metric.} Within a regime the coefficients enter the expected complete-data log-likelihood through the quadratic form $-\tfrac12\beta' M\beta + \beta' v$, with
\begin{equation}\label{eq:sm:Mv}
  M = \sum_{t=1}^{T} X_t' \Sigma_e^{-1} X_t ,
  \qquad
  v = \sum_{t=1}^{T} X_t' \Sigma_e^{-1} \hat v_t ,
\end{equation}
where $\hat v_t$ is the completed observation of period $t$ net of the smoothed latent contribution. $M$ is a \emph{generalized} least squares cross-product: it carries the estimated $\Sigma_e^{-1}$, so the effective design is $\Sigma_e^{-1/2}X_t$ and not $X_t$. Standardizing the columns of $X$ to unit Euclidean norm, the default of every penalized-regression routine, therefore does not put the covariates on a common footing. What does is to set the diagonal of $M$ to one. Write $s_j = M_{jj}^{1/2}$, $S = \mathrm{diag}(s_1,\dots,s_r)$,
\begin{equation}\label{eq:sm:std}
  \tilde M = S^{-1} M S^{-1}, \qquad \tilde v = S^{-1} v, \qquad \tilde\beta = S\beta ,
\end{equation}
so that $\tilde M_{jj} = 1$. The package applies the penalty to $\tilde\beta$ and returns $\beta = S^{-1}\tilde\beta$:
\begin{equation}\label{eq:sm:ridge-std}
  \hat{\tilde\beta} = (\tilde M + \lambda D)^{-1}\tilde v
  \quad\Longleftrightarrow\quad
  \hat\beta = (M + \lambda\, S D S)^{-1} v ,
\end{equation}
with $D$ the diagonal indicator of the penalized coordinates, the intercept excluded. On the original scale the penalty is therefore $\tfrac{\lambda}{2}\sum_j D_{jj}\,M_{jj}\,\beta_j^2$, not $\tfrac{\lambda}{2}\sum_j D_{jj}\beta_j^2$ as the compact notation
\begin{equation}\label{eq:sm:ridge-orig}
  \hat\beta = (M + \lambda D)^{-1} v
\end{equation}
of the manuscript would suggest. Each coefficient is charged in proportion to the GLS information its covariate carries. The scaling is recomputed at every M-step, since $\Sigma_e$ changes between EM iterations; it is undone before the coefficients are returned, so the iterates are well defined, and what varies across iterations is only the geometry in which a fixed $\lambda$ acts.

\paragraph{What this buys.} In the orthogonal case $\tilde M = I$ the penalized coefficient is $\tilde\beta_j = \tilde v_j/(1+\lambda)$: the unpenalized estimate multiplied by $1/(1+\lambda)$. Two invariances follow and are exact.
\begin{itemize}[leftmargin=1.4em,itemsep=2pt,topsep=3pt]
  \item \emph{Units of the response.} Replacing $z$ by $a z$ sends $v \mapsto a v$ and leaves $M$, hence $S$, unchanged, so $\tilde\beta \mapsto a\tilde\beta$ and $\beta \mapsto a\beta$.
  \item \emph{Level of the error variance.} Replacing $\Sigma_e$ by $\sigma^2\Sigma_e$ sends $M \mapsto M/\sigma^2$, $v \mapsto v/\sigma^2$, $s \mapsto s/\sigma$; hence $\tilde M$ is unchanged, $\tilde v \mapsto \tilde v/\sigma$, $\tilde\beta \mapsto \tilde\beta/\sigma$, and $\beta = S^{-1}\tilde\beta$ is unchanged.
\end{itemize}
So $\lambda$ in \eqref{eq:sm:ridge-std} is a pure number: $\lambda = 1$ halves every coefficient of an orthogonal design, whatever the units of the response and whatever the error variance of the regime.

\paragraph{The option \texttt{lambda\_scale}, and why it does nothing to a ridge.} The implementation also carries a second scaling, selected by \texttt{lambda\_scale = "relative"} (the default) or \texttt{"absolute"}, which measures a penalty against $\lambda_{\mathrm{ref}} = \max_j|\tilde v_j|$, the largest partial gradient once the unpenalized coordinates have been profiled out. It exists for the $L_1$ part of an elastic net: a lasso penalty is a \emph{threshold} compared with $\tilde v$, which carries the units of the response, and $\lambda_{\mathrm{ref}}$ removes them. It is applied to the $L_1$ part \emph{only}. A ridge has no $L_1$ part, so for the estimator of the manuscript the two settings of \texttt{lambda\_scale} return identical results. We state this explicitly because the obvious implementation, which rescales the whole penalty by $\lambda_{\mathrm{ref}}$, is wrong for the ridge: it makes $\hat\beta(\lambda)$ depend on the error variance in exactly the case in which \eqref{eq:sm:std} had removed that dependence, and on a design with a collinear pair of covariates the damage is of the order of $10^{-1}$ in the coefficients.

\paragraph{The conditioning ridge.} Independently of $\lambda$, the implementation adds a small constant, the argument \texttt{regularization} (default $0.01$), to the diagonal of the linear systems it solves, to keep them well conditioned while $\Sigma_e$ is still far from its optimum. It is not part of the statistical model, and on a well-conditioned design its effect is invisible. On a collinear design it is not: it is of the order of the smallest eigenvalue of $M$ in the offending direction, and it moves the coefficients by about $10^{-2}$. Every exact comparison below is therefore run with \texttt{regularization = 0}.

\subsection{Switching off the latent process and the spatial correlation}

Two switches remove parts of the model from the estimation rather than leaving the corresponding parameters to drift towards a boundary.

\paragraph{\texttt{latent = FALSE}.} The loading matrix is set to $K = 0$, so the latent process no longer enters the measurement equation. Its parameters $G$, $\Sigma_\eta$ and $m_0$ are then not identified --- nothing in the data speaks about a process that loads on no location --- and they are held at their input values instead of being updated by the M-step. Leaving them free is not harmless: the likelihood is flat in them, and the EM algorithm spends its iterations wandering on that flat ridge.

\paragraph{\texttt{spatial = FALSE}.} The exponential correlation $\exp(-\theta h)$ is replaced by the identity, so $\Sigma_e = (\sigma^2_\epsilon + \sigma^2_\omega) I$. Only the sum of the two variance components is then identified, and the range $\theta$ has no meaning; the Newton--Raphson update of $(\theta, \sigma^2_\epsilon/\sigma^2_\omega)$ is skipped rather than run on a direction the likelihood does not constrain. The alternative, letting $\sigma^2_\omega \to 0$, is not available by construction: the internal parameterization is $\log(\sigma^2_\epsilon/\sigma^2_\omega)$, so a zero partial sill is a boundary the algorithm can approach but never reach, and in our experiments it chased that boundary for minutes on problems that otherwise converge in a second.

\paragraph{The exact special case.} With both switches off and \texttt{regularization = 0} the model is the Gaussian linear regression $z_{it} = x_{it}'\beta + e_{it}$ with independent errors, $M = \sigma^{-2}\sum_t X_t'X_t$, and \eqref{eq:sm:ridge-std} is the ridge estimator on columns standardized to unit $X'X$ diagonal. The package reproduces it to machine precision: the ordinary least squares solution at $\lambda = 0$ to $10^{-11}$, and the ridge at $\lambda = 0.3$ to $4\times10^{-14}$. With $k > 1$ the same switches turn SC-STEM into a spatially clustered linear regression with a Potts penalty on the labels, which is the natural nested reference against which the contribution of the latent process and of the spatial covariance to an application can be measured: fit the model with and without each component and compare them out of sample.

\subsection{One \texorpdfstring{$\lambda$}{lambda} for all regimes, or one strength per regime}

With $k$ regimes the penalty is applied within each, $\hat{\tilde\beta}_h = (\tilde M_h + \lambda_h D)^{-1}\tilde v_h$, and the question is how $\lambda_h$ relates to the single hyperparameter $\lambda$. The implementation offers two answers through the argument \texttt{lambda\_by}; both keep \emph{one} hyperparameter, and they differ in what that hyperparameter holds fixed across regimes.

\paragraph{\texttt{"common"} (the default): the same proportional shrinkage.} $\lambda_h = \lambda$ for every $h$. Because the metric \eqref{eq:sm:std} is computed within each regime, $\tilde M_h$ has unit diagonal whatever the size of the regime, so in the orthogonal case every regime's coefficients are multiplied by the same factor $1/(1+\lambda)$. On the original scale the penalty on coefficient $j$ of regime $h$ is $\tfrac{\lambda}{2} M_{h,jj}\beta_{hj}^2$, and since $M_{h,jj} \approx n_h T\,\overline{x^2_j}/\sigma^2_h$ grows with the number of locations $n_h$ of the regime, the implied Gaussian prior on $\beta_{hj}$ has precision proportional to the information the regime carries --- a \emph{unit-information} prior, in the sense of the $g$-prior of \citet{Zellner1986} with $g$ proportional to the sample size.

\paragraph{\texttt{"size"}: the same prior.} $\lambda_h = \lambda\,\bar n/n_h$, with $\bar n = d/k$ the average regime size. The prior precision on the original scale becomes $\lambda\,\bar n\,M_{h,jj}/n_h \approx \lambda\,\bar n\,T\,\overline{x^2_j}/\sigma^2_h$, which no longer depends on $n_h$: every regime is shrunk towards zero by the \emph{same} prior, and a regime of half the average size is shrunk proportionally twice as hard, because its data carry half the information. A penalty of this kind is sometimes described as a penalty per observation.

\paragraph{Which one tracks the optimum.} In the orthogonal case with unit error variance on the standardized scale, the ridge parameter that minimizes the mean squared error of $\tilde\beta_{hj}$ is $\lambda^\star_{hj} = 1/\tilde\beta_{hj}^2 = 1/(M_{h,jj}\beta_{hj}^2)$. With comparable coefficients across regimes this is proportional to $1/M_{h,jj}$, that is to $1/n_h$ at fixed $T$:
\begin{equation}\label{eq:sm:lamstar}
  \lambda^\star_h \;\propto\; \frac{1}{n_h} .
\end{equation}
The \texttt{"size"} rule reproduces exactly this dependence with a single hyperparameter; the \texttt{"common"} rule ignores it, and therefore under-shrinks the small regimes and over-shrinks the large ones. The difference is immaterial with balanced regimes, where $n_h = \bar n$ and the two coincide, and it grows with the imbalance. Neither makes the penalty genuinely regime-specific, which would require one hyperparameter per regime, a grid of size $|\Lambda|^k$, and a selection rule that a regime of a dozen locations cannot support; that construction is developed elsewhere and is not used in the manuscript.

%%%%%%%%%%%%%%%%%%%%%%%%%%%%%%%%%%%%%%%%%%%%%%%%%%%%%%%%%%%%%%%%%%%%%%%%%%%%%%%
\section{Which step comes first: the ridge, the penalty and the number of regimes}\label{sm:order}
%%%%%%%%%%%%%%%%%%%%%%%%%%%%%%%%%%%%%%%%%%%%%%%%%%%%%%%%%%%%%%%%%%%%%%%%%%%%%%%

With the ridge in force the procedure carries three quantities that are chosen rather than estimated: the number of regimes $k$, the strength $\phi$ of the spatial penalty on the partition, and the strength $\lambda$ of the ridge on the coefficients. A practitioner's first question is whether they can be chosen in sequence --- fix $\lambda$ on the pooled fit, then look for $(k,\phi)$, or the reverse. They cannot, and this section says exactly why and in which direction each coupling acts.

\subsection{The objective, and where each hyperparameter enters}

With $g = (g_1,\dots,g_d)$ the labels and $E$ the edges of the neighbourhood graph, the objective is
\begin{equation}\label{eq:sm:objall}
  \mathcal{Q}(g,\Psi;\,k,\phi,\lambda)
  \;=\; \underbrace{\sum_{i=1}^{d}\ell_i(g_i;\Psi_{g_i})}_{\text{fit}}
  \;+\; \underbrace{\phi\,c\!\!\sum_{(i,j)\in E}\!\!\mathbb{1}\{g_i = g_j\}}_{\text{Potts}}
  \;-\; \underbrace{\frac{1}{2}\sum_{h=1}^{k}\lambda_h\,\tilde\beta_h' D\,\tilde\beta_h}_{\text{ridge}} .
\end{equation}
$\phi$ and $\lambda$ are weights inside \eqref{eq:sm:objall} and change the estimate; $k$ is the dimension of the model. The two penalties act on different arguments --- the Potts term on the labels, the ridge on the parameters --- and that asymmetry is the whole content of what follows.

\paragraph{Proposition S1 (the ridge leaves the label step unchanged).} \emph{At fixed $\Psi$ the map $g \mapsto \mathcal{Q}(g,\Psi)$ differs from its unpenalized counterpart by a constant, so the label update
\[
  g_i \leftarrow \argmax_h \Bigl\{\ell_i(h;\Psi_h) + \phi\,c\!\!\sum_{j:(i,j)\in E}\!\!\mathbb{1}\{h = g_j\}\Bigr\}
\]
is the same with and without the ridge.}

\emph{Proof.} Moving location $i$ from regime $h$ to regime $h'$ changes neither $\beta_h$ nor $\beta_{h'}$, which the label step holds fixed, so the third term of \eqref{eq:sm:objall} does not depend on $g$ and drops out of every arg-max. (With \texttt{lambda\_by = "size"} the weights $\lambda_h$ depend on the sizes $n_h$ and hence on $g$; they are, however, evaluated at the sizes of the parameter step and held fixed through the label step, exactly as the parameters are.) $\square$

The ridge therefore reaches the partition only \emph{through} the parameters the labels are scored against. That is not a negligible channel: shrinkage moves every $\beta_h$ towards the same point, the spread of $\ell_i(\cdot)$ across regimes narrows, and with it the evidence any single location carries about its own label.

\subsection{Three couplings, in the direction they act}

\paragraph{$\lambda$ depends on $k$.} By \eqref{eq:sm:lamstar} the optimal strength in a regime of $n_h$ locations is proportional to $1/n_h$; with $k$ balanced regimes $n_h = d/k$, so
\begin{equation}\label{eq:sm:lamk}
  \lambda^\star \;\propto\; \frac{k}{d} .
\end{equation}
The optimal shrinkage \emph{grows} with $k$, linearly. A $\lambda$ calibrated on the pooled network is too small by a factor of about $k$ for the regimes it will be used in: it under-shrinks exactly where collinearity bites hardest.

\paragraph{$k$ depends on $\lambda$.} Criteria count the effective number of coefficients $\mathrm{df}_h(\lambda) = \mathrm{tr}\{\tilde M_h(\tilde M_h + \lambda_h D)^{-1}\}$, which decreases from $r$ at $\lambda = 0$ to the number of unpenalized coordinates as $\lambda \to\infty$. The price of an additional regime under the BIC,
\begin{equation}\label{eq:sm:margk}
  \Delta_k \;=\; \log(dT)\,\bigl\{\mathrm{df}(\lambda) + q\bigr\},
\end{equation}
with $q$ the number of parameters of a regime other than the coefficients, is therefore \emph{decreasing} in $\lambda$: shrinkage makes an extra regime cheaper and pushes the selected $k$ up. The two dependencies run in opposite directions, so neither can be resolved without the other.

\paragraph{$\phi$ depends on $\lambda$, through its scale.} The default scale of the Potts term is $c = \mathrm{med}_i\{\max_h\ell_i(h) - \min_h\ell_i(h)\}/\bar m$, the typical gain from picking the best-fitting regime, per neighbour (Section~4.4 of the manuscript). Since shrinkage narrows the spread of the scores, $c$ is decreasing in $\lambda$: a \emph{fixed} $\phi$ buys less spatial smoothing at a larger $\lambda$. Two points of a grid indexed by $(\phi,\lambda)$ are comparable in $\phi$ only if $c$ is held fixed across them, which is what \texttt{phi\_scale = "raw"} does; under the default each fit calibrates its own $c$, and a joint grid must be read with that in mind.

\subsection{What follows for the search}

\begin{description}[leftmargin=1.2em,itemsep=3pt,topsep=3pt]
  \item[\textbf{Joint.}] Traverse the whole grid $(k,\phi,\lambda)$ and select on it. Correct; the cost is the product of the three grid sizes.
  \item[\textbf{Profiled, which we recommend.}] For each $(k,\phi)$ select $\lambda$; then compare the $(k,\phi)$ configurations, each at its own $\hat\lambda(k,\phi)$. The cost is that of the joint search, but the rule is applied hierarchically, which matters because $\phi$ is not chosen by a criterion at all: its Potts normalizing constant $Z(\phi,k)$ is intractable (Section~\ref{sm:ic:pen}), so $\phi$ is settled by an out-of-sample or stability argument on the partition, while $\lambda$ can be chosen by a criterion with effective degrees of freedom or by cross-validation. The $\lambda$ surface is flat relative to the $k$ surface, so a coarse logarithmic grid suffices.
  \item[\textbf{Sequential, and why not.}] Fixing $\lambda$ on the pooled fit and only then searching $(k,\phi)$ is cheap and wrong in the specific direction of \eqref{eq:sm:lamk}: it applies the shrinkage appropriate to $d$ locations inside regimes of $d/k$. The reverse order, fixing $(k,\phi)$ unpenalized and then choosing $\lambda$, is wrong in the direction of \eqref{eq:sm:margk}: the unpenalized fit charges every regime its full $r$ coefficients and so selects too few regimes for the $\lambda$ eventually used.
\end{description}

\subsection{Which likelihood the criterion sees}

It is tempting to select on the objective \eqref{eq:sm:objall} \emph{including} the ridge term. That double-counts: the reduction in flexibility the ridge produces is already carried by $\mathrm{df}(\lambda)$, and subtracting $\tfrac12\lambda\tilde\beta'D\tilde\beta$ from the log-likelihood as well charges for it twice, once through the objective and once through the dimension. The construction consistent with the effective-degrees-of-freedom literature is
\begin{equation}\label{eq:sm:icfinal}
  \mathrm{IC} \;=\; -2\,\hat\ell\bigl(\hat\Psi(\lambda)\bigr)
  \;+\; \mathrm{pen}(dT)\Bigl\{\textstyle\sum_{h=1}^k \mathrm{df}_h(\lambda) + kq\Bigr\},
\end{equation}
with $\hat\ell$ the \emph{unpenalized} log-likelihood evaluated at the \emph{penalized} estimate. The coherent alternative is Bayesian: read the ridge as a Gaussian prior and select on the marginal likelihood, which for a Gaussian prior is available in closed form. The Potts term admits the same reading only up to $Z(\phi,k)$, which is not available, so a procedure that marginalized over $\beta$ while conditioning on $\phi$ would apply two different logics to two penalties of the same objective. We do not pursue it.

%%%%%%%%%%%%%%%%%%%%%%%%%%%%%%%%%%%%%%%%%%%%%%%%%%%%%%%%%%%%%%%%%%%%%%%%%%%%%%%
\section{The data-generating process}\label{supp:dgp}\label{sm:designs}
%%%%%%%%%%%%%%%%%%%%%%%%%%%%%%%%%%%%%%%%%%%%%%%%%%%%%%%%%%%%%%%%%%%%%%%%%%%%%%%

Section~3.1 of the manuscript states the generator in equations. This section gives it as an algorithm, records the regime sizes the design produces, and states the constants.

\subsection{Constants}

\begin{center}
\begin{tabular}{lll}
\toprule
symbol & value & meaning \\
\midrule
$\nu_{\mathrm{tot}}$ & $0.4 + 2/3$ & total variance of each coordinate, held fixed \\
$\nu_{sp}(K,\omega)$ & $\nu_{\mathrm{tot}} - c_K \omega^2$ & within-cluster variance, the remainder \\
$c_K$        & $1/2,\,2/3,\,1$ & between-centre variance at $\omega=1$, for $K = 2, 3, 4$ \\
$n_{\min}$   & $6$          & smallest regime the design admits \\
$a$          & $0.7$        & temporal persistence of the covariate \\
$R_x$        & $4$          & practical range of the covariate innovations, the whole network \\
$\bar v_y$   & $0.30$       & stationary variance of every latent process \\
$\beta$      & $(2, 0.55)$  & baseline level and effect of the covariate \\
$\sigma^2_\epsilon,\ \sigma^2_\omega$ & $0.20,\ 0.20$ & baseline nugget and partial sill \\
$R$          & $2$          & baseline practical range, $\theta = 3/R$ \\
$G$          & $0.8$        & baseline persistence of the latent process \\
\bottomrule
\end{tabular}
\end{center}

The lower bound $n_{\min}$ is not the one the software enforces, which is $r+2$; it is set by the \emph{spatial} information a regime needs, since a regime estimates a range from its own pairwise distances and six locations already give fifteen of them. Below that a failure to recover the partition would be a failure to fit, not a failure to separate, and the two would be indistinguishable in the results.

\subsection{The algorithm}

\begin{algorithm}[H]
\caption{One replication of the design}\label{alg:dgp}
\begin{algorithmic}[1]
\Require $n$, $T$, $K$, overlap $\omega$, balance, scenario $s$, level $\ell$, coupling $\rho$, replication $r$
\State $(n_1,\dots,n_K) \gets \textsc{Sizes}(n, K, \text{balance})$; \; $g \gets$ labels in blocks
\State $\mu \gets \textsc{Centres}(K, \omega)$;\quad $\nu_{sp} \gets \nu_{\mathrm{tot}} - \mathrm{Var}(\mu)$
\For{$i = 1,\dots,n$} \; $s_i \gets \mu_{g_i} + N_2(0, \nu_{sp} I_2)$ \EndFor
\State $h \gets$ Euclidean distances between the $s_i$
\State $C_x \gets \exp(-3h/R_x)$; \; $L_x \gets \mathrm{chol}(C_x)$
\State $x_1 \gets L_x^\top N(0,I)$
\For{$t = 2,\dots,T$} \; $x_t \gets a\,x_{t-1} + \sqrt{1-a^2}\, L_x^\top N(0,I)$ \EndFor
\State $x \gets (x - \bar x)/\mathrm{sd}(x)$
\State $\Psi_1,\dots,\Psi_K \gets \textsc{Parameters}(s, \ell, K)$ \Comment{regime $g$ at fraction $(g-1)/(K-1)$}
\State $\Sigma_\eta \gets \Delta R_\rho \Delta$; \; $L_\eta \gets \mathrm{chol}(\Sigma_\eta)$
\State $y_1 \gets \sqrt{\bar v_y}\,\mathrm{chol}(R_\rho)^\top N(0,I)$
\For{$t = 2,\dots,T$} \; $y_t \gets \mathrm{diag}(G)\, y_{t-1} + L_\eta^\top N(0,I)$ \EndFor
\If{the regimes share $(\sigma^2_\epsilon, \sigma^2_\omega, \theta)$ and not \textsc{ForceBlock}}
  \State $L \gets \mathrm{chol}(\Sigma)$; \; \textbf{for} $t$ \textbf{do} $e_t \gets L^\top N(0,I)$
\Else
  \State \textbf{for} $g = 1,\dots,K$: $L_g \gets \mathrm{chol}(\Sigma_g)$; \textbf{for} $t$: $e_{t,I_g} \gets L_g^\top N(0,I)$
\EndIf
\For{$t = 1,\dots,T$ and $i = 1,\dots,n$}
  \State $z_{ti} \gets \beta_{0,g_i} + \beta_{1,g_i} x_{ti} + y_{t,g_i} + e_{ti}$
\EndFor
\Ensure $z$ ($T \times n$), $X = [\mathbf{1}, \mathrm{vec}(x)]$, coordinates, labels $g$, $\Psi_1,\dots,\Psi_K$
\end{algorithmic}
\end{algorithm}

Seeds are derived from the replication index alone, so a cell reproduces on its own and the same replication uses the same geometry across scenarios: differences between scenarios are then differences in the parameters and not in the draw.

\subsection{Regime sizes}

The balanced allocation is equal up to the remainder; the unbalanced one is proportional to $1:2:\cdots:K$, corrected so that no regime falls below $n_{\min}$. A cell is dropped when that correction flattens the allocation to a ratio below $1.5$ between the largest and the smallest regime, since it would then duplicate the balanced cell under another name.

\begin{center}
\begin{tabular}{lllllll}
\toprule
& & $n = 20$ & $50$ & $100$ & $200$ & $400$ \\
\midrule
balanced   & $K=3$ & 7/7/6 & 17/17/16 & 34/33/33 & 67/67/66 & 134/133/133 \\
unbalanced & $K=3$ & --- & 8/17/25 & 17/33/50 & 33/67/100 & 67/133/200 \\
\bottomrule
\end{tabular}
\end{center}

\subsection{Replications and the fitted grid}

Every replication fits the pooled model as a benchmark and the SC-STEM model over the grid $k = 1,\dots,4$ and $\phi \in \{0, 0.5, 1\}$, so that the model-selection behaviour is measured on the same replications as the estimation behaviour rather than on a sub-study. Both the pooled fit and the fit at the true number of regimes are members of that grid, so neither costs an additional run. The number of replications is $M = 100$ in every cell of the design.

Seeds are derived from the replication index, so that a cell is reproducible on its own and replication $r$ carries the same geometry and the same covariate across scenarios: a contrast between two scenarios is then a contrast between their parameters and not between two independent draws of everything else.

%%%%%%%%%%%%%%%%%%%%%%%%%%%%%%%%%%%%%%%%%%%%%%%%%%%%%%%%%%%%%%%%%%%%%%%%%%%%%%%
\section{Recovery of the partition, scenario by scenario}\label{sm:scenarios}
%%%%%%%%%%%%%%%%%%%%%%%%%%%%%%%%%%%%%%%%%%%%%%%%%%%%%%%%%%%%%%%%%%%%%%%%%%%%%%%
\TODO{For each scenario S0 to S5 and each level: the distribution of the Adjusted Rand Index across the overlap grid and the penalty grid, the bias and RMSE of the regime-specific parameters after label alignment, and the predictive accuracy relative to the pooled fit. Particular attention to S2, where the assignment score sees the dynamics only through the smoothed state, and to S3, which is the hardest case for that score and the one measuring how much of the separation the exact likelihood of the parameter step recovers on its own.}

%%%%%%%%%%%%%%%%%%%%%%%%%%%%%%%%%%%%%%%%%%%%%%%%%%%%%%%%%%%%%%%%%%%%%%%%%%%%%%%
\section{The two channels of separation}\label{sm:channels}
%%%%%%%%%%%%%%%%%%%%%%%%%%%%%%%%%%%%%%%%%%%%%%%%%%%%%%%%%%%%%%%%%%%%%%%%%%%%%%%
\TODO{The Adjusted Rand Index as a surface over the overlap and the parametric contrast, at fixed $n$ and $T$: how much spatial separation compensates for how little parametric separation, and where the procedure fails on both counts. The reference cell S5 quantifies how much of the recovery is attributable to the latent split rather than to the parameters.}

%%%%%%%%%%%%%%%%%%%%%%%%%%%%%%%%%%%%%%%%%%%%%%%%%%%%%%%%%%%%%%%%%%%%%%%%%%%%%%%
\section{Behaviour as the dimensions grow}\label{sm:dims}
%%%%%%%%%%%%%%%%%%%%%%%%%%%%%%%%%%%%%%%%%%%%%%%%%%%%%%%%%%%%%%%%%%%%%%%%%%%%%%%
\TODO{The sub-design over the whole range of $n$ and $T$: the recovery of the partition and the frequency with which the two-step rule selects the true $k$, read against the theoretical account of Section~\ref{sm:ic}. The pooled cells $K = 1$ give the false-positive rate of the rule as the dimensions grow, which is the quantity that account predicts.}

%%%%%%%%%%%%%%%%%%%%%%%%%%%%%%%%%%%%%%%%%%%%%%%%%%%%%%%%%%%%%%%%%%%%%%%%%%%%%%%
\section{Balanced and unbalanced regimes}\label{sm:balance}
%%%%%%%%%%%%%%%%%%%%%%%%%%%%%%%%%%%%%%%%%%%%%%%%%%%%%%%%%%%%%%%%%%%%%%%%%%%%%%%
\TODO{Whether an unbalanced allocation degrades the recovery of the partition beyond what the size of the smallest regime alone accounts for, and whether the minimum-size constraint of the assignment step is what binds.}

%%%%%%%%%%%%%%%%%%%%%%%%%%%%%%%%%%%%%%%%%%%%%%%%%%%%%%%%%%%%%%%%%%%%%%%%%%%%%%%
\section{Validation of the parametric bootstrap}\label{sm:boot}
%%%%%%%%%%%%%%%%%%%%%%%%%%%%%%%%%%%%%%%%%%%%%%%%%%%%%%%%%%%%%%%%%%%%%%%%%%%%%%%
\TODO{Empirical coverage of the percentile intervals for the regime-specific parameters, the ratio between the average bootstrap standard error and the Monte Carlo standard deviation of the same estimators, and the share of usable draws.}

%%%%%%%%%%%%%%%%%%%%%%%%%%%%%%%%%%%%%%%%%%%%%%%%%%%%%%%%%%%%%%%%%%%%%%%%%%%%%%%
\section{Additional material for the application}\label{sm:app}
%%%%%%%%%%%%%%%%%%%%%%%%%%%%%%%%%%%%%%%%%%%%%%%%%%%%%%%%%%%%%%%%%%%%%%%%%%%%%%%
\TODO{The full grid of fits with their admissibility diagnostics, the bootstrap co-clustering matrix, the sensitivity of the selected partition to the number of neighbors of the penalty graph, and the profile of the estimated regimes against covariates kept out of the model.}

%%%%%%%%%%%%%%%%%%%%%%%%%%%%%%%%%%%%%%%%%%%%%%%%%%%%%%%%%%%%%%%%%%%%%%%%%%%%%%%
\section{Reproduction code}\label{sm:code}
%%%%%%%%%%%%%%%%%%%%%%%%%%%%%%%%%%%%%%%%%%%%%%%%%%%%%%%%%%%%%%%%%%%%%%%%%%%%%%%

Everything in the paper is produced by the \Stem{} package \citep{StemPackage} and by a set of standalone scripts, the replication material, which is distributed separately from the package and needs nothing but an installed \Stem{}. Each script runs from whatever folder it sits in, writes its results beside itself, checks that the installed \Stem{} carries the functions it uses, and installs or updates it otherwise.

\begin{description}[leftmargin=1.6em,itemsep=3pt,topsep=3pt]
  \item[\texttt{run-simulations.R}] runs the simulation study of Section~3 of the manuscript. It carries the generator of Section~\ref{supp:dgp} and, near its top, the setup of the design: the levels of every factor, the reference cell and the blocks. \texttt{--dry} prices a run without fitting anything, calibrating the prior timings on the replications already recorded; \texttt{--coverage} runs the bootstrap experiment of Section~\ref{sm:boot} instead of the Monte Carlo. Results are appended and keyed on (cell, replication), so an interrupted run keeps what it has produced, a resumed one skips it, and blocks of replications run on different machines compose.
  \item[\texttt{analyse-simulations.R}] draws the design and turns the results into the tables and figures of the manuscript and of this document. It takes the generator and the design from the runner beside it, so the two cannot disagree.
  \item[\texttt{run-application.R}] runs the application of Section~5 of the manuscript, in stages that cache their results.
\end{description}

Every routine that uses randomness takes an explicit seed and restores the state of the random number generator on exit, so the experiments reproduce exactly. The package and the replication material are available at \url{https://github.com/PaoloMaranzano/Stem}.

\bibliographystyle{plainnat}
\bibliography{biblio}

\end{document}
